% Part: incompleteness
% Chapter: representability-in-q
% Section: introduction

\documentclass[../../../include/open-logic-section]{subfiles}

\begin{document}

\olfileid{inc}{req}{int}
\olsection{پېژندګلو}

د ناتکميلۍ قضيې پر هغو تيوريو پلي کېږي چې د محاسبه کېدونکو تابعو
په اړه بنسټيز حقايق پکې بيان او ثابتېدے شي۔ موږ به يوه ډېره لږه
داسې تيوري بيان کړو چې «$\Th{Q}$» (يا کله کله د رافايل رابنسن په
نوم «د رابنسن~$Q$») بلل کېږي۔ وبه وايو چې يوه تابعه په~$\Th{Q}$
کښې
\emph{تمثيلېدونکې} وي څه معنا لري، او بيا به لاندې خبره ثابته کړو:
\begin{quote}
  يوه تابعه په~$\Th{Q}$ کښې هله او يوازې هله تمثيلېدونکې ده چې
  محاسبه کېدونکې وي۔
\end{quote}
دا له يوې خوا د محاسبې يو بل مدل راکوي۔ خو له بلې خوا به ئې د دې د
ښودلو دپاره هم وکاروو چې سټ~$\Setabs{!A}{\Th{Q} \Proves !A}$ د
پرېکړې وړ نۀ دے؛ د درېدنې مسئله به ورته راکمه کړو۔ تر هغه وخته چې
بحث بشپړوو، تر دې به ډېرې پياوړې پايلې مو ثابتې کړې وي۔

د~$\Th{Q}$ ژبه د حساب ژبه ده؛ $\Th{Q}$ له لاندې بديهي اصولو څخه
جوړه ده (چې د !!{identity} لرونکي لومړۍ درجې منطق له نورو بديهي
اصولو او قاعدو سره يوځای کارول کېږي):
\begin{align*}
& \lforall[x][\lforall[y][(\eq[x'][y'] \lif \eq[x][y])]] \tag{$!Q_1$}\\
& \lforall[x][\eq/[\Obj 0][x']] \tag{$!Q_2$}\\
& \lforall[x][(\eq[x][\Obj 0] \lor \lexists[y][\eq[x][y']])] \tag{$!Q_3$}\\
& \lforall[x][\eq[(x + \Obj 0)][x]] \tag{$!Q_4$}\\
& \lforall[x][\lforall[y][\eq[(x + y')][(x + y)']]] \tag{$!Q_5$}\\
& \lforall[x][\eq[(x \times \Obj 0)][\Obj 0]] \tag{$!Q_6$}\\
& \lforall[x][\lforall[y][\eq[(x \times y')][((x \times y) + x)]]] \tag{$!Q_7$}\\
& \lforall[x][\lforall[y][(x < y \liff \lexists[z][\eq[(z' + x)][y]])]] \tag{$!Q_8$}
\end{align*}
د هر طبيعي عدد~$n$ دپاره عددنښه~$\num{n}$ داسې تعريف کړئ چې ترم
$\Obj{0}^{\prime\prime\ldots\prime}$ وي او د پريم نښو ټول شمېر ئې~$n$ وي۔ نو
$\num{0}$ يوازې !!{constant}~$\Obj{0}$ دے، $\num{1}$،
$\Obj{0}'$ دے، $\num{2}$، $\Obj{0}''$ دے، او همداسې وړاندې۔

د حساب د يوې تيورۍ په توګه~$\Th{Q}$ \emph{بېخي} کمزورې ده؛ د
بېلګې په توګه ان د~$\lforall[x][\eq/[x][x']]$ يا
$\lforall[x][\lforall[y][(x + y) = (y + x)]]$ په شان ډېر ساده
حقايق هم پکې نۀ شئ ثابتولے۔ خو وبه وينو چې د~$\Th{Q}$ د پام وړتيا
تر ډېره \emph{له همدې} کمزورۍ څخه راځي۔ په حقيقت کښې دومره پياوړې
ده چې د ناتکميلۍ قضيه پرې پلې شي، خو له دې زياته نۀ ده۔ د~$\Th{Q}$
د پام وړتيا بل لامل دا دے چې د بديهي اصولو سټ ئې \emph{متناهي} دے۔

له~$\Th{Q}$ څخه يوه پياوړې تيوري (چې د \emph{پيانو حساب}
$\Th{PA}$ بلل کېږي) په~$\Th{Q}$ کښې د استقرا شېما په زياتولو
ترلاسه کېږي:
\[
(!A(\Obj 0) \land \lforall[x][(!A(x) \lif !A(x'))]) \lif \lforall[x][!A(x)]
\]
دلته~$!A(x)$ هر فارمول دے۔ که~$!A(x)$ له~$x$ پرته نور ازاد
!!{variable}s ولري، په سر کښې کلي کمیت ايښوونکي ورزياتوو چې ټول وتړي
(څو د استقرا د شېما اړونده بېلګه !!a{sentence} وي)۔ د بېلګې په
توګه، که~$!A(x, y)$ کښې !!{variable}~$y$ هم ازاد وي، اړونده بېلګه
دا ده:
\[
\lforall[y][((!A(\Obj 0) \land \lforall[x][(!A(x) \lif !A(x'))]) \lif
  \lforall[x][!A(x)])]
\]
د استقرا د شېما د بېلګو په کارولو سره د~$\Th{PA}$ له بديهي اصولو
څخه د~$\Th{Q}$ په پرتله ډېر څه ثابتېدے شي۔ په حقيقت کښې د طبيعي
عددونو په اړه د داسې «طبيعي» بيانونو موندل ډېر کار غواړي چې د پيانو
په حساب کښې نۀ ثابتېږي!{}

\begin{defn}
\ollabel{defn:representable-fn}
  له طبيعي عددونو څخه طبيعي عددونو ته تابعه
  $f(x_0,\ldots,x_k)$ هله \emph{په~$\Th{Q}$ کښې تمثيلېدونکې}
  بلل کېږي چې داسې فارمول~$!A_f(x_0,\dots,x_k,y)$ وي چې هر
  کله~$f(n_0,\dots,n_k) = m$ وي، $\Th{Q}$ لاندې دواړه ثابت کړي:
\begin{enumerate}
\item\ollabel{defn:rep:a} $!A_f(\num{n_0}, \dots, \num{n_k}, \num{m})$
\item\ollabel{defn:rep:b} $\lforall[y][(!A_f(\num{n_0}, \dots,
\num{n_k}, y) \lif \num{m} = y)]$۔
\end{enumerate}
\end{defn}

د تعريف د بيان نورې برابرې لارې هم شته؛ د بېلګې په توګه په برابر
ډول دا غوښتلے شو چې~$\Th{Q}$،
$\lforall[y][(!A_f(\num{n_0}, \dots,
\num{n_k}, y) \liff \eq[y][\num{m}])]$ ثابت کړي۔

\begin{thm}
\ollabel{thm:representable-iff-comp}
يوه تابعه په~$\Th{Q}$ کښې هله او يوازې هله تمثيلېدونکې ده چې
محاسبه کېدونکې وي۔
\end{thm}

د قضيې ثبوت دوه لوري لري۔ له کيڼ څخه ښي لور ته ثبوت، کله چې د نحو
حسابي کول جوړ شي، نسبتاً ساده دے۔ بل لور زيات کار غواړي۔ بنسټيزه
مفکوره دا ده: «عمومي بازګشتي» د «محاسبه کېدونکي» د کره کولو يوه لار
ټاکو، او ښيو چې هره عمومي بازګشتي تابعه په~$\Th{Q}$ کښې
تمثيلېدونکې ده۔ په ياد ولرئ، تابعه هله عمومي بازګشتي ده چې له
$\Zero$، تالي تابع~$\Succ$، او پروجکشن تابعو~$\Proj{n}{i}$ څخه د
ترکيب، بنسټيز بازګښت، او منظم اقل موندلو په کارولو سره تعريفېدے شي۔
نو د دې د ښودلو يوه لار چې هره عمومي بازګشتي تابعه په~$\Th{Q}$ کښې
تمثيلېدونکې ده، دا ده چې وښيو بنسټيزې تابعې تمثيلېدونکې دي، او هر
کله چې ځينې تابعې تمثيلېدونکې وي، د هغو له ترکيب، بنسټيز بازګښت، او
منظم اقل موندلو څخه تعريف شوې تابعې هم تمثيلېدونکې وي۔ په بله وينا،
ښودلے شو چې بنسټيزې تابعې تمثيلېدونکې دي، او تمثيلېدونکې تابعې د
ترکيب، بنسټيز بازګښت، او منظم اقل موندلو «لاندې تړلې» دي۔ دا خبره
باور راکوي چې هره عمومي بازګشتي تابعه تمثيلېدونکې ده۔

خو دا ګام چې وښيو تمثيلېدونکې تابعې د بنسټيز بازګښت لاندې تړلې دي،
ستونزمن راوځي۔ له دې ګام څخه د ډډې دپاره لومړے ښيو چې په حقيقت کښې
له بنسټيز بازګښت پرته کار کولے شو۔ يعنې ښيو چې هره عمومي بازګشتي
تابعه يوازې د بنسټيزو تابعو، ترکيب، او منظم اقل موندلو په کارولو
تعريفېدے شي۔ د دې دپاره ښيو چې بنسټيز بازګښت په رښتيا د يوه ځانګړي
منظم اقل موندلو په وسيله ترسره کېدے شي۔ خو د دې د سم کار دپاره بايد
يو شمېر نورې بنسټيزې تابعې هم ورزياتې کړو: جمع، ضرب، او د عينيت د
اړيکې ځانګړونکې تابعه~$\Char{=}$۔ بيا قضيه په دې ښودلو ثابتولے شو
چې دا \emph{ټولې} بنسټيزې تابعې په~$\Th{Q}$ کښې تمثيلېدونکې دي، او
تمثيلېدونکې تابعې د ترکيب او منظم اقل موندلو لاندې تړلې دي۔

\end{document}
